% Generated draft: requires mathematical and source review.
\documentclass[11pt]{article}
\usepackage[margin=25mm]{geometry}
\usepackage{fontspec}
\setmainfont{Latin Modern Roman}
\usepackage{amsmath,amssymb,mathtools,unicode-math}
\setmathfont{Latin Modern Math}
\usepackage{enumitem,needspace,xurl,hyperref}
\setlength{\parindent}{0pt}
\setlength{\parskip}{5pt}
\setlength{\emergencystretch}{3em}
\newcommand{\N}{\mathbb{N}}
\newcommand{\Z}{\mathbb{Z}}
\newcommand{\Q}{\mathbb{Q}}
\newcommand{\R}{\mathbb{R}}
\newcommand{\C}{\mathbb{C}}
\newcommand{\F}{\mathbb{F}}
\newcommand{\A}{\mathbb{A}}
\newcommand{\PP}{\mathbb{P}}
\newcommand{\E}{\mathbb{E}}
\newcommand{\eps}{\varepsilon}
\newcommand{\dd}{\,\mathrm{d}}
\begin{document}
\section*{UnsolvedMath ID: 30004480}
\textbf{Source problem code:} catalogue entry
\subsection{Definitions and mathematical statement}
The supplied record uses \(\mathrm{KV}\) to denote a symmetry group of solutions of the Kashiwara--Vergne problem and refers to a prounipotent Grothendieck--Teichm\"uller group. It does not specify the versions of these groups, their base field, the solution space, or the symmetry action. These descriptions therefore identify the intended objects only provisionally.

\textbf{Source-level assertion.} The symmetry group \(\mathrm{KV}\) of Kashiwara--Vergne solutions is conjecturally closely related to the prounipotent Grothendieck--Teichm\"uller group.

\textbf{Unresolved formulation.} The phrase ``closely related'' specifies no mathematical relation, map, equality, isomorphism, or quantified property. According to the supplied source-verification note, the report describes an approach illuminating a relationship conjectured elsewhere, without stating a precise standalone conjecture. The record is marked \texttt{statement\_status: unrecoverable}, and no clean statement is supplied. Consequently, the evidence does not determine a precise mathematical proposition; the assertion above is retained for provenance and requires recovery of the intended conjecture before formalization.

The supplied literature assessment also describes an extracted welded-foam claim, but no welded-foam objects or precise claim are given. Its statements that the Kashiwara--Vergne problem itself is solved and that its symmetry group is related to Grothendieck--Teichm\"uller groups are supplied background assessments, not independently verified conclusions of this draft.
\subsection{Short English statement}
The record suggests a conjectural relationship between Kashiwara--Vergne symmetries and a prounipotent Grothendieck--Teichm\"uller group, but does not specify what that relationship means. It cannot presently be turned into a precise standalone conjecture. The supplied open label concerns this imprecise extracted relationship, while the supplied assessment describes the Kashiwara--Vergne problem itself as solved.
\subsection{Sources}
\begin{itemize}
\item Local UnsolvedMath registry record, permanent ID \(30004480\), source problem code \texttt{catalogue entry}, title ``Imprecise Kashiwara--Vergne Symmetry Relation,'' in the Oberwolfach Reports Open Problems collection. The record gives proposed year \(2020\) and no named proposer.
\item Supplied report citation: \emph{Topologie} (2021), Oberwolfach Reports, DOI 10.4171/owr/2020/21, pp.~1083--1137. \url{https://doi.org/10.4171/owr/2020/21}. The supplied background describes Oberwolfach Reports as documenting workshops and their open-problem discussions.
\item Dancso--Halacheva--Robertson, \emph{A topological characterisation of the KV groups}. Supplied literature source: \url{https://arxiv.org/abs/2106.02373}.
\item Dancso--Hogan--Robertson, \emph{A knot-theoretic approach to KV and GT groups}. Supplied literature source: \url{https://arxiv.org/abs/2211.11370}.
\item OWR Topology report. Supplied literature source: \url{https://ems.press/content/serial-article-files/46858}.
\item Literature assessment dated exactly \texttt{2026-08-22}: supplied dated curation, not independently verified. It labels the extracted problem open, describes the Kashiwara--Vergne problem itself as solved, and characterizes the extracted welded-foam relationship as conjectural, without a proved equality or complete formulation. The supplied background requests independent verification of time-sensitive status and cited claims.
\item The supplied statement-verification note says that the official report does not formulate a precise standalone claim and that the original record is retained publicly for provenance. The conjecture said to occur elsewhere is not identified precisely in the supplied evidence.
\end{itemize}
\end{document}
